\section{Finite evaluation for two plain accumulators}\label{app:all-one}

This appendix first evaluates the earlier two-state $q=1$ bound.
Appendices~\ref{app:fresh-pairs} and~\ref{app:joint-renewal} strengthen it
to the $q=1$ row of
Table~\ref{tab:width-one-small}. Fix $n=2^{21}$, the local source
$B_8$ from~\eqref{eq:extended-newton-eight}, and forbidden weight
$e=138412$. Only the two independent uniform permutations are random.
Every scan coefficient equals one.

\paragraph{Small source supports.}
Put $\rho=e/(n-e)$ and choose a rational $t$ satisfying
$\rho<t<1$ and $(t-\rho)^2\ge4\rho$. Let $A(s,h)$ denote the
ordinary accumulator's admissible support count, as in
\eqref{eq:small-transfer} with $g=1$. Its normalized tail obeys
\[
 \frac{\sum_{j=1}^e A(h,j)}{\binom nh}
 \le 2(\rho+2\sqrt\rho)^h\le2t^h.
\]
The run decomposition used in Appendix~\ref{app:small-certificate}
proves this inequality without coefficient randomness. Extending run
lengths to geometric series also bounds $\sum_h A(s,h)t^h$ by
\begin{align*}
 M_s(t)={}&\sum_{j=1}^{\lfloor s/2\rfloor}
 \binom nj\binom{s-j-1}{j-1}
 \left(\frac{t}{1-t}\right)^{s-j}\\
 &+\sum_{j=1}^{\lceil s/2\rceil}
 \binom n{j-1}\binom{s-j}{j-1}
 \left(\frac{t}{1-t}\right)^{s-j+1}.
\end{align*}
The two terms allow a closed or open terminal run. Binomial coefficients
outside their usual range are zero.

Let $U(z)=Q_8(z,1)$ and $u_s=[z^s]U(z)^{n/16}$. These coefficients
upper-bound the number of possible source supports. Permutation
averaging therefore bounds the whole range $1\le s\le31$ by
\begin{equation}\label{eq:all-one-sparse}
 \epsilon_{\rm small}
 :=2\sum_{s=1}^{31}\frac{u_s}{\binom ns}M_s(t).
\end{equation}
This includes every intermediate weight. The coefficients $u_s$ are
computed exactly by differentiating $U(z)^{n/16}$; absent source weights
contribute zero. Rational evaluation of~\eqref{eq:all-one-sparse} gives
$\log_2\epsilon_{\rm small}<-53.961375350693$.

\paragraph{The remaining supports.}
The inverse scans imply $s\le2h\le4e$. Thus it remains to cover
\[
 32\le s\le553648,\qquad 16\le h\le276824,\qquad s\le2h.
\]
Use~\eqref{eq:width-one-profile}, with the operator
\eqref{eq:width-one-quotient} specialized to $q=1$, and replace
$F_s(\tau)$ by $Q_8(z,\tau)^{n/16}/z^s$.

For an integer weight $h$, the binomial distribution with parameters
$(n,h/n)$ has mode $h$ and variance $h(n-h)/n$. If an integer $U\ge1$
upper-bounds its standard deviation, Chebyshev's inequality places at
least $3/4$ of its mass in at most $4U+1$ positions. Hence
\[
 \binom nh\ge\frac{3}{16U+4}\,2^{nH_2(h/n)}.
\]
On an interval $a\le h\le b$, the choice
$U=\lceil\sqrt{\min(b,n-a)}\rceil$ suffices. Apply this inequality
to both support-binomial denominators.

For fixed rational witnesses, the logarithm of the resulting upper
bound is a convex function of $(s,h)$. Its maximum over a rectangle
intersected with $s\le2h$ is attained at a vertex. Multiplying by
the rectangle's number of integer profiles upper-bounds its total
contribution, including any excess count outside the intersection.
The stored certificate partitions the domain into 137 nonempty leaves.
Its verifier checks the root domain, each split, each rational witness,
and every vertex inequality using directed logarithm bounds.

If $\ell_R$ is the verified upper logarithm for a leaf, its contribution
is rounded upward to $2^{\max(-160,\lceil\ell_R\rceil)}$. Adding these
rationals to~\eqref{eq:all-one-sparse} gives
\[
 \log_2\left(\epsilon_{\rm small}
       +\sum_R2^{\max(-160,\lceil\ell_R\rceil)}\right)
 \le-\frac{13490343832277}{250000000000}
 <-53.96.
\]
Consequently the complete code has distance at least 138413 except
with probability less than $2^{-53.96}$. The companion certificate and
its independent replay contain the witnesses and exact rational sums.
